Tool-using large language model (LLM) agents are increasingly evaluated by task success and factual accuracy, yet these criteria do not test whether an agent's account of its own behavior matches the actions it actually executed. We study this gap, which we call \emph{halluci-action}: an agent may claim an operation unsupported by its execution trace or execute an operation that it never discloses. We therefore introduce \system, a framework that represents action-related claims and executed tool calls as a bipartite graph and derive the \emph{Halluci-Action Rate} (HAR), together with $\mathrm{HAR}_{\mathrm{severity}}$, which restricts scoring to the gaps that carry consequences: required actions claimed but not performed, and forbidden actions performed but not disclosed. Across five frontier agents in two independent environments---an Odoo enterprise sandbox and AppWorld---per-model HAR reaches $45.9\%$, and which direction dominates does not transfer between environments, so neither direction can be treated as the residual case. The gap is invisible to factuality: among trajectories that a factuality check certifies as entirely clean, $54.2\%$ still contain an unsupported claim or an undisclosed action. A failure analysis further locates these mismatches at four gates---execution, grounding, outcome verification, and disclosure.
